\section{Problem Formulation}
\label{sec:formulation}

Let $u$ denote a user, $x$ a current request, $d$ its domain, $t$ its task type,
$m$ a model provider/snapshot, and $P_u$ the portable preference profile. A
selector $S(P_u,d,t,x)$ returns a bounded set of relevant preference records,
while compiler $C$ maps those records to an instruction context $c$. Generation
is
\begin{equation}
  y \sim p_m\!\left(y \mid x, C(S(P_u,d,t,x))\right).
\end{equation}
The current request is a hard precedence constraint: if $x$ conflicts with a
stored default, $x$ wins. Selection is therefore not merely nearest-neighbor
retrieval; it combines scope compatibility, evidence authority, confidence,
recency, locks, and conflict resolution.

For subjective utility $Q(u,x,y)$, the primary estimand is the average
within-user contrast between domain-dynamic personalization and no
personalization over a prespecified population of users, tasks, and model
families. Portability additionally requires that $P_u$ is unchanged across $m$;
thin message-format adapters do not count as model-specific preference storage.
We therefore distinguish \emph{syntactic portability} (the profile can be parsed
and compiled for a new model) from \emph{behavioral conformance} (the resulting
response satisfies the same applicable policy without increasing correctness,
safety, privacy, or cost failures beyond a prespecified margin).

The design treats correctness, safety, privacy exposure, context cost, and
preference adherence as distinct outcomes. A system that increases stylistic
adherence while decreasing correctness does not unconditionally improve quality.
